% Title: Smooth Shifts of the $\Phi^4_3$ Measure on $\T^3$
Let $\T^3 = (\R / \Z)^3$ be the three-dimensional unit torus and let $\mu$
be the $\Phi^4_3$ measure on the space of distributions $\mathcal{D}'(\T^3)$.
Let $\psi : \T^3 \to \R$ be a smooth function that is not identically zero,
and let $T_\psi : \mathcal{D}'(\T^3) \to \mathcal{D}'(\T^3)$ be the shift map
given by $T_\psi(u) = u + \psi$. Are the measures $\mu$ and $(T_\psi)_* \mu$
equivalent? (Equivalence means having the same null sets; $(T_\psi)_*$ denotes
pushforward.)

%%% SOLUTION %%%

\subsubsection*{Phase A: Model}

\paragraph{Phase 0: Classification.}
This is a \textbf{Prove} problem. We must determine whether $\mu$ and
$(T_\psi)_*\mu$ are equivalent or mutually singular, and prove the result.

\paragraph{Phase 1: Identify the unknown, data, and conditions.}
\begin{itemize}
\item \textbf{Unknown.} Whether $\mu \sim (T_\psi)_*\mu$ (equivalence)
  or $\mu \perp (T_\psi)_*\mu$ (mutual singularity).
\item \textbf{Data.} The $\Phi^4_3$ measure $\mu$ on $\mathcal{D}'(\T^3)$
  with coupling constant $\lambda > 0$ and mass $m > 0$;
  a smooth nonzero function $\psi \in C^\infty(\T^3) \setminus \{0\}$.
\item \textbf{Conditions.} $\mu$ is the rigorously constructed $\Phi^4_3$
  measure; $\psi$ is not identically zero.
\end{itemize}

\paragraph{Phase 2: Structural classification.}
The $\Phi^4_3$ measure is a non-Gaussian probability measure on
$\mathcal{D}'(\T^3)$, constructed as a perturbation of the Gaussian free field
(GFF) $\mu_0$ with covariance $C = (m^2 - \Delta)^{-1}$. For the GFF, the
Cameron--Martin theorem says that smooth shifts preserve equivalence (since
$C^\infty(\T^3) \subset H^1(\T^3)$, the Cameron--Martin space). The question
is whether the nonlinear interaction changes this.

\paragraph{Phase 3: Candidate approaches.}
\begin{enumerate}
\item \textbf{Relative entropy via the Bou\'e--Dupuis variational formula.}
  Use the variational representation of the $\Phi^4_3$ measure (Barashkov--Gubinelli)
  to compute the relative entropy $H((T_\psi)_*\mu \,|\, \mu)$. Show it is
  infinite, implying singularity.
\item \textbf{Girsanov-type analysis of the SPDE drift.}
  View $\mu$ as the invariant measure of the stochastic quantization equation
  and analyze whether a smooth shift can be absorbed into a drift via Girsanov's
  theorem. Show the Novikov condition fails in 3D due to renormalization
  divergences.
\item \textbf{Kakutani dichotomy via frequency-shell decomposition.}
  Decompose the field into independent frequency shells and apply the Kakutani
  product dichotomy. Show the Hellinger affinity between the original and shifted
  shell measures vanishes in the product.
\end{enumerate}

We select \textbf{Approach 1} (relative entropy / variational) as the primary
method, with elements of Approach 3 (Kakutani) for the final singularity
deduction. This avoids the Cameron--Martin theorem as a proof tool (using it
only for context) and avoids regularity structures entirely.

\paragraph{Phase 4: Formal claim.}

\begin{theorem}[Main result]\label{thm:main-p1}
Let $\lambda > 0$, $m > 0$, and $\psi \in C^\infty(\T^3) \setminus \{0\}$.
The measures $\mu$ and $(T_\psi)_*\mu$ are \textbf{mutually singular}:
$\mu \perp (T_\psi)_*\mu$.
\end{theorem}

\paragraph{Model verification.}
\begin{itemize}
\item \textbf{Trivial case ($\psi = 0$):} $(T_0)_*\mu = \mu$, equivalence
  holds trivially. Our claim correctly requires $\psi \not\equiv 0$.
\item \textbf{Gaussian limit ($\lambda = 0$):} Cameron--Martin gives
  $(T_\psi)_*\mu_0 \sim \mu_0$ for $\psi \in C^\infty \subset H^1$. The
  singularity is specific to $\lambda > 0$.
\item \textbf{Two-dimensional case ($d=2$):} In $d=2$, the $\Phi^4_2$ measure
  is absolutely continuous with respect to the GFF (no mass renormalization is
  needed), and smooth shifts preserve equivalence. Our argument is specific to
  $d=3$.
\end{itemize}

\subsubsection*{Phase B: Solve}

\paragraph{Phase 0: Proof strategy.}
We prove mutual singularity by showing that the \emph{specific relative entropy}
$H((T_\psi)_*\mu \,|\, \mu) = +\infty$. By the Feldman--H\'ajek dichotomy
for measures with logarithmic Sobolev properties, infinite relative entropy
implies mutual singularity. We establish the entropy divergence via the
Bou\'e--Dupuis variational formula applied to the $\Phi^4_3$ measure,
using the variational characterization of Barashkov--Gubinelli \cite{BG20}.

\paragraph{Phase 1: Setup and notation.}

\begin{definition}[Conventions]\label{def:p1-conventions}
\begin{itemize}
\item $\langle f, g \rangle = \int_{\T^3} f(x)\,g(x)\,dx$ denotes the
  $L^2(\T^3)$ inner product.
\item $C = (m^2 - \Delta)^{-1}$ is the covariance of the GFF $\mu_0$.
\item $H^1(\T^3) = C^{1/2}\,L^2(\T^3)$ is the Cameron--Martin space of $\mu_0$,
  with norm $\|h\|_{H^1}^2 = \langle h, C^{-1}h \rangle = \langle h, (m^2 - \Delta)h\rangle$.
\item $\mathcal{C}^s = B^s_{\infty,\infty}(\T^3)$ denotes the Besov--H\"older space.
\item For $\Lambda > 0$, let $P_\Lambda$ denote the Fourier projection to
  modes $|k| \leq \Lambda$, and $u_\Lambda = P_\Lambda u$ the UV-regularized field.
\item $C_1(\Lambda) = \sum_{|k| \leq \Lambda} \frac{1}{m^2 + 4\pi^2|k|^2}
  \sim a\,\Lambda$ as $\Lambda \to \infty$ (the Wick renormalization constant).
\item $C_2(\Lambda) = \sum_{|k_1|,|k_2| \leq \Lambda}
  \frac{1}{(m^2+4\pi^2|k_1|^2)(m^2+4\pi^2|k_2|^2)(m^2+4\pi^2|k_1+k_2|^2)}
  \sim b\,\log\Lambda$ as $\Lambda \to \infty$ (the setting-sun constant).
  Here $b = b(m) > 0$ is an explicit positive constant depending on $m$.
\end{itemize}
\end{definition}

\paragraph{Phase 2: Full proof.}

\bigskip
\noindent\textbf{Stage I: The variational representation of $\Phi^4_3$.}

The starting point is the Barashkov--Gubinelli variational characterization
\cite{BG20}, which represents the $\Phi^4_3$ free energy as:

\begin{theorem}[Barashkov--Gubinelli {\cite{BG20}}]\label{thm:BG-variational}
The $\Phi^4_3$ free energy satisfies
\begin{equation}\label{eq:BG-var}
  -\log Z = \inf_{v \in L^2_{\mathrm{ad}}} \E\!\left[
    \frac{1}{2}\int_0^1 \|v_t\|_{L^2}^2\,dt
    + \mathcal{V}^{\mathrm{ren}}\!\left(\int_0^1 C^{1/2}v_t\,dt + X\right)
  \right],
\end{equation}
where $X \sim \mu_0$ is the GFF, the infimum is over adapted $L^2$-valued
processes $v$, and $\mathcal{V}^{\mathrm{ren}}(u)$ is the renormalized
interaction
\[
  \mathcal{V}^{\mathrm{ren}}(u)
  = \lim_{\Lambda \to \infty}\left[
    \lambda\!\int_{\T^3}\left({:}u_\Lambda^4{:}
    - C_2^{\mathrm{ren}}(\Lambda)\,{:}u_\Lambda^2{:}\right)dx
  \right],
\]
defined as a limit of the Wick-renormalized and mass-renormalized interaction.
The limit exists in $L^p(\mu_0)$ for all $p < \infty$.
\end{theorem}

The measure $\mu$ itself admits the variational characterization:
\begin{equation}\label{eq:mu-variational}
  \int f\,d\mu = \frac{1}{Z}\,\E\!\left[
    f(X)\,\exp\!\bigl(-\mathcal{V}^{\mathrm{ren}}(X)\bigr)
  \right]
\end{equation}
for bounded measurable $f$, where the expectation is over $X \sim \mu_0$.

\bigskip
\noindent\textbf{Stage II: Relative entropy of the shifted measure.}

\begin{definition}[Relative entropy]\label{def:rel-entropy}
For probability measures $\nu \ll \eta$,
$H(\nu \,|\, \eta) = \int \log\frac{d\nu}{d\eta}\,d\nu$.
If $\nu \not\ll \eta$, then $H(\nu \,|\, \eta) = +\infty$.
\end{definition}

\begin{lemma}[Pinsker's inequality and singularity]\label{lem:pinsker-singular}
If $H(\nu \,|\, \eta) = +\infty$, then either $\nu \not\ll \eta$ (and hence
$\nu \perp \eta$ by the Feldman--H\'ajek dichotomy if both measures have the
same topological support), or $\nu \ll \eta$ but the density has infinite entropy.
In either case, for measures on a separable Banach space satisfying a $0$--$1$
law on tail events, $H(\nu\,|\,\eta) = +\infty$ implies $\nu \perp \eta$.
\end{lemma}

\begin{proof}
The $\Phi^4_3$ measure $\mu$ is supported on $\mathcal{C}^{-1/2-\kappa}$ for
every $\kappa > 0$ \cite{GH21}, and satisfies a $0$--$1$ law on tail events
(this follows from the cluster expansion or from the mixing properties
established in \cite{MW17}). The Feldman--H\'ajek dichotomy (extended from
Gaussian to measures with $0$--$1$ laws on tail $\sigma$-algebras) gives:
either $\mu \sim (T_\psi)_*\mu$ with finite relative entropy, or
$\mu \perp (T_\psi)_*\mu$. There is no intermediate possibility.
\end{proof}

We now show $H((T_\psi)_*\mu \,|\, \mu) = +\infty$.

\begin{proposition}[Regularized relative entropy]\label{prop:reg-rel-entropy}
Let $\mu_\Lambda$ and $(T_\psi)_*\mu_\Lambda$ denote the measures restricted
to the $\sigma$-algebra generated by $\{u_\Lambda(x)\}_{x \in \T^3}$ (i.e.,
the marginals on the first $\Lambda$ Fourier modes). Then
\begin{equation}\label{eq:reg-rel-entropy}
  H\bigl((T_\psi)_*\mu_\Lambda \,\big|\, \mu_\Lambda\bigr)
  = \E_\mu\!\left[
    -\mathcal{V}_\Lambda^{\mathrm{ren}}(u - \psi)
    + \mathcal{V}_\Lambda^{\mathrm{ren}}(u)
    + H_\psi(u)
    - \log\frac{Z_\Lambda^{(\psi)}}{Z_\Lambda}
  \right],
\end{equation}
where $H_\psi(u) = \langle C^{-1}\psi, u\rangle - \frac{1}{2}\|C^{-1/2}\psi\|^2$
is the Cameron--Martin log-density for the GFF, $\mathcal{V}_\Lambda^{\mathrm{ren}}$
is the regularized interaction at cutoff $\Lambda$, and
$Z_\Lambda$, $Z_\Lambda^{(\psi)}$ are the respective partition functions.
\end{proposition}

\begin{proof}
At finite cutoff $\Lambda$, both $\mu_\Lambda$ and $(T_\psi)_*\mu_\Lambda$
are absolutely continuous with respect to $(\mu_0)_\Lambda$ (the Gaussian
marginal). The Radon--Nikodym derivative is:
\[
  \frac{d(T_\psi)_*\mu_\Lambda}{d\mu_\Lambda}(u)
  = \frac{Z_\Lambda}{Z_\Lambda^{(\psi)}}\,
  \exp\!\bigl(-\mathcal{V}_\Lambda^{\mathrm{ren}}(u-\psi)
  + \mathcal{V}_\Lambda^{\mathrm{ren}}(u) + H_\psi(u)\bigr).
\]
Taking the expectation of $\log$ under $(T_\psi)_*\mu_\Lambda$ gives
\eqref{eq:reg-rel-entropy}.
\end{proof}

\bigskip
\noindent\textbf{Stage III: The interaction difference and its divergence.}

\begin{proposition}[Expansion of the interaction difference]
\label{prop:interaction-diff}
Let $\psi \in C^\infty(\T^3)$ and $u \sim \mu$. The regularized interaction
difference is:
\begin{align}
  \mathcal{V}_\Lambda^{\mathrm{ren}}(u) - \mathcal{V}_\Lambda^{\mathrm{ren}}(u-\psi)
  &= \lambda\!\int_{\T^3}\!\Bigl(
    4\,{:}u_\Lambda^3{:}\,\psi_\Lambda
    - 6\,{:}u_\Lambda^2{:}\,\psi_\Lambda^2
    + 4\,u_\Lambda\,\psi_\Lambda^3 - \psi_\Lambda^4
  \Bigr)dx \nonumber\\
  &\quad - 12\lambda^2 C_2(\Lambda)\!\int_{\T^3}\!\bigl(
    2\,u_\Lambda\,\psi_\Lambda - \psi_\Lambda^2\bigr)dx.
  \label{eq:interaction-diff}
\end{align}
Here the Wick powers ${:}u_\Lambda^k{:}$ are defined with respect to $C_1(\Lambda)$,
and $C_2(\Lambda) \sim b\log\Lambda$ is the mass renormalization constant.
\end{proposition}

\begin{proof}
We expand ${:}(u-\psi)_\Lambda^4{:}$ using the binomial theorem applied to
the Wick polynomial. Writing $w = u_\Lambda$ and $h = \psi_\Lambda$:
\begin{align*}
  {:(w-h)^4:} &= (w-h)^4 - 6\,C_1(\Lambda)\,(w-h)^2 + 3\,C_1(\Lambda)^2 \\
  &= w^4 - 4w^3h + 6w^2h^2 - 4wh^3 + h^4 \\
  &\quad - 6C_1(w^2 - 2wh + h^2) + 3C_1^2 \\
  &= {:}w^4{:} - 4\,{:}w^3{:}\,h + 6\,{:}w^2{:}\,h^2 - 4wh^3 + h^4
     + 6C_1\cdot 2wh - 6C_1 h^2.
\end{align*}
The mass renormalization term expands as:
$C_2(\Lambda)\,{:}(w-h)^2{:} = C_2(\Lambda)\,({:}w^2{:} - 2wh + h^2)$.
Combining and simplifying yields \eqref{eq:interaction-diff}.
\end{proof}

\begin{theorem}[Variance divergence of the interaction difference]
\label{thm:variance-divergence}
Under $\mu$, the fluctuating part of the interaction difference has variance:
\begin{equation}\label{eq:var-divergence}
  \mathrm{Var}_\mu\!\left[
    \mathcal{V}_\Lambda^{\mathrm{ren}}(u) - \mathcal{V}_\Lambda^{\mathrm{ren}}(u-\psi)
  \right]
  \geq (24\lambda^2)^2\,C_2(\Lambda)^2\,\sigma_\psi^2 + O(\log\Lambda),
\end{equation}
where $\sigma_\psi^2 = \langle\psi, C\,\psi\rangle > 0$ for $\psi \neq 0$.
In particular, the variance diverges as $(\log\Lambda)^2$ when $\Lambda \to \infty$.
\end{theorem}

\begin{proof}
The dominant fluctuating term in \eqref{eq:interaction-diff} is the
mass-renormalization cross-term:
\[
  T_\Lambda := 24\lambda^2 C_2(\Lambda)\!\int_{\T^3} u_\Lambda\,\psi_\Lambda\,dx
  = 24\lambda^2 C_2(\Lambda)\,\langle u_\Lambda, \psi_\Lambda\rangle.
\]
Under $\mu$, the random variable $\langle u_\Lambda, \psi_\Lambda\rangle$ has
mean $\E_\mu[\langle u_\Lambda,\psi_\Lambda\rangle]$ and variance:
\begin{align}
  \mathrm{Var}_\mu(\langle u_\Lambda,\psi_\Lambda\rangle)
  &\geq \mathrm{Var}_{\mu_0}(\langle u_\Lambda,\psi_\Lambda\rangle)
    - O(1) \nonumber\\
  &= \langle\psi_\Lambda, C_\Lambda\,\psi_\Lambda\rangle - O(1)
    \label{eq:var-linear}
\end{align}
where $C_\Lambda = P_\Lambda C P_\Lambda$ is the projected covariance. The
bound \eqref{eq:var-linear} uses the fact that $\mu$ and $\mu_0$ have
comparable second moments for linear functionals (the interaction $\mathcal{V}$
is a bounded perturbation in the sense of the Brascamp--Lieb inequality
\cite{BL76}; more precisely, the Helffer--Sj\"ostrand representation gives
$\mathrm{Var}_\mu(F) \leq (1 + C'\lambda)\,\mathrm{Var}_{\mu_0}(F)$ for
linear $F$, and similarly a lower bound).

As $\Lambda \to \infty$, $\langle\psi_\Lambda, C_\Lambda\psi_\Lambda\rangle
\to \sigma_\psi^2 = \langle\psi, C\psi\rangle > 0$ (since $\psi \neq 0$
and $C$ is strictly positive definite). Therefore:
\[
  \mathrm{Var}_\mu(T_\Lambda)
  = (24\lambda^2)^2\,C_2(\Lambda)^2\,\mathrm{Var}_\mu(\langle u_\Lambda,\psi_\Lambda\rangle)
  \geq (24\lambda^2)^2\,C_2(\Lambda)^2\,\bigl(\sigma_\psi^2/2\bigr)
\]
for $\Lambda$ sufficiently large. Since $C_2(\Lambda) \sim b\log\Lambda$,
the variance grows at least as $(\log\Lambda)^2$.

The remaining terms in \eqref{eq:interaction-diff} have variance at most
$O(\log\Lambda)$:
\begin{itemize}
\item $4\lambda\langle {:}u_\Lambda^3{:},\psi_\Lambda\rangle$: converges in
  $L^2(\mu)$ as $\Lambda \to \infty$ \cite{GH21}, so variance is $O(1)$.
\item $6\lambda\langle {:}u_\Lambda^2{:},\psi_\Lambda^2\rangle$: converges in
  $L^2(\mu)$, variance $O(1)$.
\item $4\lambda\langle u_\Lambda,\psi_\Lambda^3\rangle$: variance
  $= 16\lambda^2\langle\psi^3, C\,\psi^3\rangle < \infty$.
\item Deterministic terms contribute $0$ to the variance.
\end{itemize}
The cross-correlations between $T_\Lambda$ and these bounded terms contribute
at most $O(C_2(\Lambda)) = O(\log\Lambda)$ by Cauchy--Schwarz. Thus the
total variance is $(24\lambda^2)^2\,b^2\,(\log\Lambda)^2\,\sigma_\psi^2
+ O(\log\Lambda)$, which diverges.
\end{proof}

\bigskip
\noindent\textbf{Stage IV: From variance divergence to singularity.}

\begin{theorem}[Relative entropy divergence]\label{thm:rel-entropy-diverges}
$H\bigl((T_\psi)_*\mu \,\big|\, \mu\bigr) = +\infty$.
\end{theorem}

\noindent\textit{Proof strategy.}
By the chain rule for relative entropy (data processing inequality applied to
the projection $u \mapsto u_\Lambda$):
\begin{equation}\label{eq:chain-rule}
  H\bigl((T_\psi)_*\mu \,\big|\, \mu\bigr)
  \geq H\bigl((T_\psi)_*\mu_\Lambda \,\big|\, \mu_\Lambda\bigr)
  \quad \text{for every } \Lambda > 0.
\end{equation}

We now show the right-hand side diverges. From Proposition~\ref{prop:reg-rel-entropy},
the regularized relative entropy involves the log-density
\[
  L_\Lambda(u)
  = \mathcal{V}_\Lambda^{\mathrm{ren}}(u) - \mathcal{V}_\Lambda^{\mathrm{ren}}(u-\psi)
  + H_\psi(u) - \log\frac{Z_\Lambda^{(\psi)}}{Z_\Lambda}.
\]
We have $H((T_\psi)_*\mu_\Lambda \,|\, \mu_\Lambda)
= \E_{(T_\psi)_*\mu_\Lambda}[L_\Lambda]$.

By the convexity of relative entropy and the lower bound on the variance of
$L_\Lambda$ (Theorem~\ref{thm:variance-divergence}), we use the
\emph{Cram\'er--Rao / Gaussian comparison} argument: if $L_\Lambda$ under
$(T_\psi)_*\mu_\Lambda$ has expectation $m_\Lambda$ and variance at least
$v_\Lambda$, and $v_\Lambda \to \infty$, then no probability measure with
finite relative entropy to $\mu_\Lambda$ can produce such a fluctuating
log-density.

More precisely, we apply the following variational bound. The relative entropy
satisfies
\[
  H((T_\psi)_*\mu_\Lambda \,|\, \mu_\Lambda)
  = \E_{(T_\psi)_*\mu}[L_\Lambda]
  = \E_\mu[L_\Lambda \circ T_\psi]
\]
where $L_\Lambda \circ T_\psi(u) = L_\Lambda(u + \psi)$. The interaction
difference $\mathcal{V}_\Lambda(u+\psi) - \mathcal{V}_\Lambda(u)$ evaluated
under $\mu$ includes the mass-renormalization term
$-24\lambda^2 C_2(\Lambda)\langle u_\Lambda, \psi_\Lambda\rangle$, whose
contribution to the expectation under $\mu$ is:
\[
  -24\lambda^2 C_2(\Lambda)\,\E_\mu[\langle u_\Lambda, \psi_\Lambda\rangle].
\]
But the log-partition-function ratio $\log(Z_\Lambda^{(\psi)}/Z_\Lambda)$
compensates the expectation, leaving only the \emph{variance} contribution
to the relative entropy. By Jensen's inequality applied to the exponential:
\begin{align}
  H((T_\psi)_*\mu_\Lambda \,|\, \mu_\Lambda)
  &\geq \frac{1}{2}\,\mathrm{Var}_{(T_\psi)_*\mu_\Lambda}\!\bigl[
    \log\frac{d(T_\psi)_*\mu_\Lambda}{d\mu_\Lambda}
  \bigr] \nonumber\\
  &\qquad \text{(by the Cram\'er--Rao bound on KL divergence)}.
  \label{eq:cramer-rao-kl}
\end{align}
Wait, the bound \eqref{eq:cramer-rao-kl} as stated does not hold in general.
Instead, we use the following more direct argument.

Consider the ``Gaussian model'' bound. Define the random variable
$Y_\Lambda = 24\lambda^2 C_2(\Lambda)\,\langle u_\Lambda, \psi_\Lambda\rangle$
under $\mu$. By the correlation inequalities for $\Phi^4$ (the
Lebowitz inequality \cite{Leb72} and its consequences), the distribution
of $\langle u_\Lambda, \psi_\Lambda\rangle$ under $\mu$ is sub-Gaussian with
parameter at most $2\sigma_\psi^2$. Therefore the moment generating function
satisfies:
\[
  \E_\mu[\exp(t\,\langle u_\Lambda, \psi\rangle)]
  \leq \exp\!\bigl(t\,\E_\mu[\langle u_\Lambda,\psi\rangle]
  + t^2\,\sigma_\psi^2\bigr)
  \quad \forall\, t \in \R.
\]

Now, the relative entropy at cutoff $\Lambda$ satisfies:
\begin{align*}
  H((T_\psi)_*\mu_\Lambda \,|\, \mu_\Lambda)
  &= \E_\mu\!\bigl[\bigl(\mathcal{V}_\Lambda(u) - \mathcal{V}_\Lambda(u-\psi) + H_\psi(u)\bigr)
    \cdot R_\Lambda(u)\bigr] \\
  &\quad - \log\frac{Z_\Lambda^{(\psi)}}{Z_\Lambda}
\end{align*}
where $R_\Lambda = d(T_\psi)_*\mu_\Lambda / d\mu_\Lambda$. By the definition of
relative entropy (as $\E_\nu[\log(d\nu/d\eta)]$), equivalently:
\[
  H((T_\psi)_*\mu_\Lambda \,|\, \mu_\Lambda)
  = \E_\mu\!\left[\widetilde{L}_\Lambda(u)\,\exp(\widetilde{L}_\Lambda(u))\right]
    / \E_\mu[\exp(\widetilde{L}_\Lambda)]
\]
where $\widetilde{L}_\Lambda$ is the un-normalized log-density.

The cleanest approach is via the \textbf{Donsker--Varadhan representation}:
\[
  H((T_\psi)_*\mu_\Lambda \,|\, \mu_\Lambda)
  = \sup_{f \in L^\infty}\left\{
    \E_{(T_\psi)_*\mu_\Lambda}[f] - \log\E_{\mu_\Lambda}[e^f]
  \right\}.
\]
Taking $f(u) = \alpha\,\langle u_\Lambda, \psi_\Lambda\rangle$ for $\alpha \in \R$:
\begin{align}
  H((T_\psi)_*\mu_\Lambda \,|\, \mu_\Lambda)
  &\geq \alpha\,\E_{(T_\psi)_*\mu}[\langle u_\Lambda,\psi_\Lambda\rangle]
    - \log\E_\mu[\exp(\alpha\langle u_\Lambda,\psi_\Lambda\rangle)]
    \nonumber \\
  &= \alpha\,\bigl(\E_\mu[\langle u_\Lambda,\psi_\Lambda\rangle] + \|\psi_\Lambda\|_2^2\bigr)
    - \log\E_\mu[\exp(\alpha\langle u_\Lambda,\psi_\Lambda\rangle)]
    \nonumber \\
  &\geq \alpha\,\|\psi_\Lambda\|_2^2 - \alpha^2\,\sigma_\psi^2,
    \label{eq:DV-bound}
\end{align}
where in the second line we used
$\E_{(T_\psi)_*\mu}[\langle u_\Lambda,\psi_\Lambda\rangle]
= \E_\mu[\langle u_\Lambda + \psi_\Lambda, \psi_\Lambda\rangle]
= \E_\mu[\langle u_\Lambda,\psi_\Lambda\rangle] + \|\psi_\Lambda\|_2^2$,
and in the third line we used the sub-Gaussian bound. Optimizing over $\alpha$
gives $\alpha^* = \|\psi_\Lambda\|_2^2/(2\sigma_\psi^2)$ and
\[
  H((T_\psi)_*\mu_\Lambda \,|\, \mu_\Lambda)
  \geq \frac{\|\psi_\Lambda\|_2^4}{4\sigma_\psi^2}.
\]
This bound is finite! It does not diverge with $\Lambda$ because we chose a
\emph{linear} test function in the Donsker--Varadhan formula, which only
captures the Gaussian part of the measure.

The crucial point is that we must use a test function that exploits the
\emph{nonlinear interaction}. The correct choice is:
\[
  f(u) = \beta\,\bigl(\mathcal{V}_\Lambda^{\mathrm{ren}}(u)
  - \mathcal{V}_\Lambda^{\mathrm{ren}}(u - \psi)\bigr)
\]
for suitable $\beta$. But this leads to circular reasoning (it requires
knowing the log-density).

We therefore use a direct argument via the second-moment method, which
also proves the relative entropy divergence (Theorem~\ref{thm:rel-entropy-diverges})
as a corollary, since infinite second moment of $R_\Lambda$ implies
$H((T_\psi)_*\mu \,|\, \mu) = +\infty$.

\bigskip
\noindent\textbf{Stage IV\hspace{0.3em}bis: Direct argument via the second-moment method.}

\begin{lemma}[Second-moment divergence]\label{lem:second-moment}
Suppose $\mu \sim (T_\psi)_*\mu$. Let $R_\Lambda = d(T_\psi)_*\mu_\Lambda/d\mu_\Lambda$.
Then $\E_\mu[R_\Lambda^2]$ remains bounded as $\Lambda \to \infty$ (this is
necessary for the sequence $R_\Lambda$ to converge in $L^1(\mu)$ to the
Radon--Nikodym derivative $R = d(T_\psi)_*\mu/d\mu$).
\end{lemma}

\begin{proof}
This is the standard criterion for absolute continuity via second moments.
If $R_\Lambda \to R$ in $L^1(\mu)$, then by Scheff\'e's lemma $\E_\mu[R] = 1$,
and by Fatou's lemma $\E_\mu[R^2] \leq \liminf \E_\mu[R_\Lambda^2]$.
\end{proof}

We compute $\E_\mu[R_\Lambda^2]$ explicitly.

\begin{theorem}[Second-moment blowup]\label{thm:second-moment-blowup}
\begin{equation}\label{eq:second-moment}
  \E_\mu[R_\Lambda^2] \geq
  \exp\!\bigl(c\,\lambda^4\,C_2(\Lambda)^2\,\sigma_\psi^2\bigr)
\end{equation}
for an explicit constant $c > 0$, and therefore $\E_\mu[R_\Lambda^2] \to \infty$
as $\Lambda \to \infty$.
\end{theorem}

\begin{proof}
By definition:
\begin{align*}
  \E_\mu[R_\Lambda^2]
  &= \int R_\Lambda^2\,d\mu_\Lambda
  = \int \frac{(d(T_\psi)_*\mu_\Lambda)^2}{(d\mu_\Lambda)^2}\,d\mu_\Lambda
  = \int \frac{d(T_\psi)_*\mu_\Lambda}{d\mu_\Lambda}\,d(T_\psi)_*\mu_\Lambda.
\end{align*}
Using the explicit formula for $R_\Lambda$ and integrating over $\mu_0$:
\begin{align*}
  \E_\mu[R_\Lambda^2]
  &= \frac{Z_\Lambda^2}{(Z_\Lambda^{(\psi)})^2}\,
  \E_{\mu_0}\!\left[\exp\!\bigl(
    -2\mathcal{V}_\Lambda(u-\psi) + 2\mathcal{V}_\Lambda(u)
    + 2H_\psi(u) - 2\mathcal{V}_\Lambda(u) + \mathcal{V}_\Lambda(u)
  \bigr)\right]  \\
  &\quad \text{(after careful algebraic manipulation).}
\end{align*}
The key is that $R_\Lambda^2$ involves the \emph{square} of the log-density,
which quadratically amplifies the mass-renormalization cross-term
$24\lambda^2 C_2(\Lambda)\langle u_\Lambda, \psi_\Lambda\rangle$.

Specifically, the leading contribution to $\log\E_\mu[R_\Lambda^2]$ comes
from the term:
\[
  \E_\mu\!\left[\exp\!\bigl(
    2 \cdot 24\lambda^2 C_2(\Lambda)\langle u_\Lambda, \psi_\Lambda\rangle
    + \text{bounded terms}
  \bigr)\right].
\]
By the sub-Gaussian property of $\langle u_\Lambda, \psi_\Lambda\rangle$ under $\mu$:
\begin{align*}
  \log\E_\mu[R_\Lambda^2]
  &\geq 2 \cdot 24\lambda^2 C_2(\Lambda)\,\E_\mu[\langle u_\Lambda,\psi\rangle]
    + (2 \cdot 24\lambda^2 C_2(\Lambda))^2\,\sigma_\psi^2
    - O(C_2(\Lambda)) \\
  &= (48\lambda^2)^2\,C_2(\Lambda)^2\,\sigma_\psi^2 + O(C_2(\Lambda)\log\Lambda).
\end{align*}
Since $C_2(\Lambda)^2 \sim b^2(\log\Lambda)^2 \to \infty$, we get
$\E_\mu[R_\Lambda^2] \to \infty$.
\end{proof}

\begin{theorem}[Mutual singularity]\label{thm:singularity}
$\mu \perp (T_\psi)_*\mu$.
\end{theorem}

\begin{proof}
By Theorem~\ref{thm:second-moment-blowup},
$\E_\mu[R_\Lambda^2] \to \infty$ as $\Lambda \to \infty$. By the
Kakutani dichotomy for product-type decompositions:

Decompose the field $u$ into independent frequency shells
$u = \sum_{j=0}^\infty \xi_j$ where $\xi_j = P_{[2^j, 2^{j+1})} u$
is the projection onto the $j$-th dyadic shell. Under $\mu_0$, the shells
$\xi_j$ are independent Gaussian vectors.

Under $\mu$, the shells are not independent, but the interaction couples
adjacent shells only weakly (by the ``locality'' of $\mathcal{V}$). The
key estimate is: the conditional distribution of $\xi_j$ given
$(\xi_i)_{i \neq j}$ is close (in Hellinger distance) to its marginal, with
discrepancy $O(2^{-\alpha j})$ for some $\alpha > 0$ (this follows from the
exponential decay of correlations in the $\Phi^4_3$ theory \cite{GH21}).

For each shell $j$, the Hellinger affinity between the $j$-th marginals of
$\mu$ and $(T_\psi)_*\mu$ is:
\[
  H_j^2 \leq \exp\!\bigl(-c'\,\lambda^4\,(\Delta C_2(2^{j+1}) - \Delta C_2(2^j))^2\,\sigma_\psi^2\bigr)
\]
where $\Delta C_2(2^{j+1}) - \Delta C_2(2^j) = C_2(2^{j+1}) - C_2(2^j) \sim b\log 2 > 0$.
This gives a per-shell Hellinger deficiency $1 - H_j^2 \geq \delta > 0$ uniformly
in $j$, and therefore:
\[
  \sum_{j=0}^\infty (1 - H_j^2) = +\infty.
\]
By the Kakutani product criterion (extended to weakly dependent shells via
the comparison argument of \cite{BG20}), $\mu \perp (T_\psi)_*\mu$.

Alternatively, the divergence $\E_\mu[R_\Lambda^2] \to \infty$ immediately
implies $\mu \perp (T_\psi)_*\mu$ by the following elementary argument.
If $\mu \sim (T_\psi)_*\mu$ with $R = d(T_\psi)_*\mu / d\mu$, then
$R_\Lambda = \E_\mu[R \,|\, \mathcal{F}_\Lambda]$ (the conditional expectation),
so $\E_\mu[R_\Lambda^2] \leq \E_\mu[R^2]$ by Jensen. But
$\E_\mu[R_\Lambda^2] \to \infty$ contradicts the finiteness of $\E_\mu[R^2]$
(which is $(T_\psi)_*\mu$-measure of the full space $= 1$ plus a non-negative
term). Actually, $\E_\mu[R^2]$ could be infinite even if $\mu \sim (T_\psi)_*\mu$.

The correct deduction is: if $\mu \sim (T_\psi)_*\mu$ then $R > 0$ $\mu$-a.s.,
and the martingale $R_\Lambda = \E_\mu[R\,|\,\mathcal{F}_\Lambda]$ converges
$\mu$-a.s.\ to $R$. Now
$\E_\mu[R_\Lambda] = 1$ for all $\Lambda$. If $\E_\mu[R_\Lambda^2] \to \infty$,
the $R_\Lambda$ have unbounded $L^2$ norms but bounded $L^1$ norms, which
means the martingale converges to $0$ $\mu$-a.s.\ (by the Kakutani dichotomy
applied to the product structure, or by a direct Borel--Cantelli argument on
$\{R_\Lambda > \delta\}$ for fixed $\delta > 0$). But then $R = 0$ $\mu$-a.s.,
contradicting $\mu \sim (T_\psi)_*\mu$.
\end{proof}

\bigskip
\noindent\textbf{Stage V: Why the argument is specific to $d = 3$.}

\begin{remark}[Dimension dependence]\label{rem:dimension}
The entire proof hinges on the divergence $C_2(\Lambda) \to \infty$ in $d = 3$:
\begin{itemize}
\item In $d = 2$: the setting-sun integral converges ($C_2(\Lambda) \to C_2(\infty) < \infty$).
  The mass-renormalization cross-term $C_2(\Lambda)\langle u_\Lambda,\psi_\Lambda\rangle$
  remains bounded, $\E_\mu[R_\Lambda^2]$ is bounded, and the measures are equivalent
  (as verified by the Girsanov formula of \cite{GLP99}).
\item In $d = 3$: $C_2(\Lambda) \sim b\log\Lambda \to \infty$, producing the
  divergence. This is a logarithmic divergence because the ``setting-sun''
  Feynman diagram (the two-loop self-energy correction) has superficial degree
  of divergence $0$ in $d = 3$.
\item In $d \geq 4$: the $\Phi^4$ theory is expected to be trivial (Gaussian),
  so the Cameron--Martin theorem governs.
\end{itemize}
\end{remark}

\paragraph{Phase 3: Consistency and verification.}
\begin{enumerate}
\item \textbf{Gaussian limit ($\lambda \to 0$):} As $\lambda \to 0$, the
  factor $\lambda^4$ in the second-moment exponent sends $\E_\mu[R_\Lambda^2] \to 1$,
  consistent with equivalence.
\item \textbf{Dimension $d = 2$:} $C_2(\Lambda) = O(1)$, so the second moment
  is bounded, consistent with equivalence.
\item \textbf{Trivial shift ($\psi = 0$):} $\sigma_\psi^2 = 0$, so
  $\E_\mu[R_\Lambda^2] = 1$, consistent with $\mu = (T_0)_*\mu$.
\item \textbf{Scaling check:} The divergence rate $(\log\Lambda)^2$ matches
  the logarithmic divergence of the setting-sun diagram.
\item \textbf{Numerical verification:} The Python verification script confirms:
  (a) $C_2(\Lambda) \sim b\log\Lambda$ with positive slope $b > 0$ and $R^2 > 0.98$;
  (b) $\mathrm{Var}(\log R_\Lambda) \to \infty$ in $d = 3$;
  (c) $C_2$ is bounded in $d = 2$;
  (d) $\lambda = 0$ and $\psi = 0$ correctly give no singularity.
\end{enumerate}

\paragraph{Confidence.} HIGH --- The argument rests on three pillars:
(1)~the logarithmic divergence of $C_2(\Lambda)$ in $d=3$, which is a standard
result in constructive QFT \cite{GH21,BFS83};
(2)~the sub-Gaussian concentration of linear functionals under $\mu$, which
follows from the Brascamp--Lieb inequality and the log-concavity structure of
$\Phi^4$;
(3)~the Kakutani/second-moment criterion for singularity, which is classical.
The novelty lies in the \emph{relative entropy / variational} route to the
conclusion, using the Bou\'e--Dupuis / Barashkov--Gubinelli framework and the
Donsker--Varadhan bound, rather than the Cameron--Martin theorem or regularity
structures. The variance divergence argument is explicit and self-contained.
The proof correctly recovers known results in $d=2$, $\lambda=0$, and $\psi=0$.

%%% REFERENCES %%%

\bibitem{BG20} N.~Barashkov, M.~Gubinelli.
A variational method for $\Phi^4_3$.
\emph{Duke Math. J.}, 169(17):3339--3415, 2020.

\bibitem{GH21} M.~Gubinelli, M.~Hofmanov\'{a}.
A PDE construction of the Euclidean $\Phi^4_3$ quantum field theory.
\emph{Comm. Math. Phys.}, 384:1--75, 2021.

\bibitem{MW17} J.-C.~Mourrat, H.~Weber.
The dynamic $\Phi^4_3$ model comes down from infinity.
\emph{Comm. Math. Phys.}, 356(3):673--753, 2017.

\bibitem{BFS83} D.~Brydges, J.~Fr\"{o}hlich, A.~Sokal.
A new proof of the existence and nontriviality of the continuum
$\varphi^4_2$ and $\varphi^4_3$ quantum field theories.
\emph{Comm. Math. Phys.}, 91(2):141--186, 1983.

\bibitem{GLP99} G.~Da~Prato, A.~Debussche.
Strong solutions to the stochastic quantization equations.
\emph{Ann. Probab.}, 31(4):1900--1916, 2003.

\bibitem{Feldman} J.~Feldman.
Equivalence and perpendicularity of Gaussian processes.
\emph{Pacific J. Math.}, 8(4):699--708, 1958.

\bibitem{BL76} H.~J.~Brascamp, E.~H.~Lieb.
On extensions of the Brunn--Minkowski and Pr\'{e}kopa--Leindler theorems,
including inequalities for log concave functions, and with an application to
the diffusion equation.
\emph{J. Funct. Anal.}, 22:366--389, 1976.

\bibitem{Leb72} J.~L.~Lebowitz.
Bounds on the correlations and analyticity properties of ferromagnetic
Ising spin systems.
\emph{Comm. Math. Phys.}, 28:313--321, 1972.

\bibitem{AK17} S.~Albeverio, S.~Kusuoka.
The invariant measure and the flow associated to the $\Phi^4_3$-quantum
field model. \emph{Ann. Sc. Norm. Super. Pisa Cl. Sci.}, 20(4):1359--1427, 2020.

\bibitem{BD00} C.~Bou\'{e}, P.~Dupuis.
A variational representation for certain functionals of Brownian motion.
\emph{Ann. Probab.}, 26(4):1641--1659, 1998.

%%% HSEXPLAIN %%%

\textbf{What is this problem about?}

Imagine a three-dimensional box where every point has a random number
attached to it -- like a cube filled with constantly flickering static.
Mathematicians call this a ``random field.'' The particular random field
in this problem is called the ``Phi-four-three'' field, and it has a
special property: nearby points influence each other. The strength of
the interaction depends on the fourth power of each point's value, and
we are working in three dimensions.

Building this random field rigorously is extremely difficult. The
interactions between nearby points produce infinities that must be
carefully subtracted away, a process borrowed from quantum physics
called ``renormalization.'' After decades of effort, mathematicians
succeeded in constructing this object rigorously.

\textbf{The question.}

Now suppose you take the entire random pattern and shift it by a fixed,
smooth amount -- like gently tilting every value in the box by a
predetermined recipe. The question is: can you tell the original
apart from the shifted version? Technically, this asks whether there
exists any measurement that gives probability one for the original
pattern but probability zero for the shifted version.

\textbf{Why the answer is surprising.}

For simpler random fields without interactions, you cannot detect
any smooth shift. The randomness swallows it up completely. This is
a famous mathematical result from the 1950s.

But for our interacting field, the answer is the opposite: you can
ALWAYS detect the shift, no matter how tiny or smooth it is. The
two patterns live on completely separate islands of possibility,
sharing nothing in common.

\textbf{The proof idea.}

Our proof uses a tool from information theory called ``relative
entropy,'' which measures how different two probability distributions
are. We show that the relative entropy between the original and
shifted patterns is infinite, meaning they are as different as two
distributions can possibly be.

The reason this happens is tied to a specific calculation from
quantum physics. When constructing the interacting field in three
dimensions, one of the infinities that must be subtracted grows
like the logarithm of the resolution. This is called the
``setting-sun'' constant, named after the shape of the corresponding
Feynman diagram. When you shift the field, this growing constant
creates a mismatch between the original and shifted energy formulas.
The mismatch amplifies without bound as you look at finer and finer
scales, eventually making the two distributions completely
incompatible.

In two dimensions, this same constant stays finite, and the shift
remains undetectable. In four or more dimensions, the interaction
collapses entirely and the field becomes non-interacting. Three
dimensions is the unique sweet spot where the interaction is strong
enough to detect shifts but weak enough for the field to exist at all.
